% Options for packages loaded elsewhere
% Options for packages loaded elsewhere
\PassOptionsToPackage{unicode}{hyperref}
\PassOptionsToPackage{hyphens}{url}
%
\documentclass[
  ignorenonframetext,
  aspectratio=169,
]{beamer}
\newif\ifbibliography
\usepackage{pgfpages}
\setbeamertemplate{caption}[numbered]
\setbeamertemplate{caption label separator}{: }
\setbeamercolor{caption name}{fg=normal text.fg}
\beamertemplatenavigationsymbolsempty
% remove section numbering
\setbeamertemplate{part page}{
  \centering
  \begin{beamercolorbox}[sep=16pt,center]{part title}
    \usebeamerfont{part title}\insertpart\par
  \end{beamercolorbox}
}
\setbeamertemplate{section page}{
  \centering
  \begin{beamercolorbox}[sep=12pt,center]{section title}
    \usebeamerfont{section title}\insertsection\par
  \end{beamercolorbox}
}
\setbeamertemplate{subsection page}{
  \centering
  \begin{beamercolorbox}[sep=8pt,center]{subsection title}
    \usebeamerfont{subsection title}\insertsubsection\par
  \end{beamercolorbox}
}
% Prevent slide breaks in the middle of a paragraph
\widowpenalties 1 10000
\raggedbottom
\AtBeginPart{
  \frame{\partpage}
}
\AtBeginSection{
  \ifbibliography
  \else
    \frame{\sectionpage}
  \fi
}
\AtBeginSubsection{
  \frame{\subsectionpage}
}
\usepackage{iftex}
\ifPDFTeX
  \usepackage[T1]{fontenc}
  \usepackage[utf8]{inputenc}
  \usepackage{textcomp} % provide euro and other symbols
\else % if luatex or xetex
  \usepackage{unicode-math} % this also loads fontspec
  \defaultfontfeatures{Scale=MatchLowercase}
  \defaultfontfeatures[\rmfamily]{Ligatures=TeX,Scale=1}
\fi
\usepackage{lmodern}

\usetheme[]{Boadilla}
\usecolortheme[]{default}
\usefonttheme[]{professionalfonts}
\ifPDFTeX\else
  % xetex/luatex font selection
\fi
% Use upquote if available, for straight quotes in verbatim environments
\IfFileExists{upquote.sty}{\usepackage{upquote}}{}
\IfFileExists{microtype.sty}{% use microtype if available
  \usepackage[]{microtype}
  \UseMicrotypeSet[protrusion]{basicmath} % disable protrusion for tt fonts
}{}
\makeatletter
\@ifundefined{KOMAClassName}{% if non-KOMA class
  \IfFileExists{parskip.sty}{%
    \usepackage{parskip}
  }{% else
    \setlength{\parindent}{0pt}
    \setlength{\parskip}{6pt plus 2pt minus 1pt}}
}{% if KOMA class
  \KOMAoptions{parskip=half}}
\makeatother


\usepackage{longtable,booktabs,array}
\newcounter{none} % for unnumbered tables
\usepackage{calc} % for calculating minipage widths
\usepackage{caption}
% Make caption package work with longtable
\makeatletter
\def\fnum@table{\tablename~\thetable}
\makeatother
\usepackage{graphicx}
\makeatletter
\newsavebox\pandoc@box
\newcommand*\pandocbounded[1]{% scales image to fit in text height/width
  \sbox\pandoc@box{#1}%
  \Gscale@div\@tempa{\textheight}{\dimexpr\ht\pandoc@box+\dp\pandoc@box\relax}%
  \Gscale@div\@tempb{\linewidth}{\wd\pandoc@box}%
  \ifdim\@tempb\p@<\@tempa\p@\let\@tempa\@tempb\fi% select the smaller of both
  \ifdim\@tempa\p@<\p@\scalebox{\@tempa}{\usebox\pandoc@box}%
  \else\usebox{\pandoc@box}%
  \fi%
}
% Set default figure placement to htbp
\def\fps@figure{htbp}
\makeatother





\setlength{\emergencystretch}{3em} % prevent overfull lines

\providecommand{\tightlist}{%
  \setlength{\itemsep}{0pt}\setlength{\parskip}{0pt}}



 


\makeatletter
\@ifpackageloaded{tcolorbox}{}{\usepackage[skins,breakable]{tcolorbox}}
\@ifpackageloaded{fontawesome5}{}{\usepackage{fontawesome5}}
\definecolor{quarto-callout-color}{HTML}{909090}
\definecolor{quarto-callout-note-color}{HTML}{0758E5}
\definecolor{quarto-callout-important-color}{HTML}{CC1914}
\definecolor{quarto-callout-warning-color}{HTML}{EB9113}
\definecolor{quarto-callout-tip-color}{HTML}{00A047}
\definecolor{quarto-callout-caution-color}{HTML}{FC5300}
\definecolor{quarto-callout-color-frame}{HTML}{acacac}
\definecolor{quarto-callout-note-color-frame}{HTML}{4582ec}
\definecolor{quarto-callout-important-color-frame}{HTML}{d9534f}
\definecolor{quarto-callout-warning-color-frame}{HTML}{f0ad4e}
\definecolor{quarto-callout-tip-color-frame}{HTML}{02b875}
\definecolor{quarto-callout-caution-color-frame}{HTML}{fd7e14}
\makeatother
\makeatletter
\@ifpackageloaded{caption}{}{\usepackage{caption}}
\AtBeginDocument{%
\ifdefined\contentsname
  \renewcommand*\contentsname{Table of contents}
\else
  \newcommand\contentsname{Table of contents}
\fi
\ifdefined\listfigurename
  \renewcommand*\listfigurename{List of Figures}
\else
  \newcommand\listfigurename{List of Figures}
\fi
\ifdefined\listtablename
  \renewcommand*\listtablename{List of Tables}
\else
  \newcommand\listtablename{List of Tables}
\fi
\ifdefined\figurename
  \renewcommand*\figurename{Figure}
\else
  \newcommand\figurename{Figure}
\fi
\ifdefined\tablename
  \renewcommand*\tablename{Table}
\else
  \newcommand\tablename{Table}
\fi
}
\@ifpackageloaded{float}{}{\usepackage{float}}
\floatstyle{ruled}
\@ifundefined{c@chapter}{\newfloat{codelisting}{h}{lop}}{\newfloat{codelisting}{h}{lop}[chapter]}
\floatname{codelisting}{Listing}
\newcommand*\listoflistings{\listof{codelisting}{List of Listings}}
\makeatother
\makeatletter
\makeatother
\makeatletter
\@ifpackageloaded{caption}{}{\usepackage{caption}}
\@ifpackageloaded{subcaption}{}{\usepackage{subcaption}}
\makeatother

\usepackage{bookmark}
\IfFileExists{xurl.sty}{\usepackage{xurl}}{} % add URL line breaks if available
\urlstyle{same}
\hypersetup{
  pdftitle={Chapter 2: Risk Sharing --- Reinsurance and Deductibles},
  hidelinks,
  pdfcreator={LaTeX via pandoc}}


\title{\texorpdfstring{Chapter 2: Risk Sharing --- Reinsurance and
Deductibles}{Chapter 2: Risk Sharing --- Reinsurance and Deductibles}}
\author{}
\date{}

\begin{document}
\frame{\titlepage}


\section{Introduction to Risk
Sharing}\label{introduction-to-risk-sharing}

\begin{frame}{Risk Sharing in Insurance}
\protect\phantomsection\label{risk-sharing-in-insurance}
Insurance is a form of \textbf{risk sharing}.

When a loss occurs, the financial burden may be shared among:

\begin{itemize}
\tightlist
\item
  the \textbf{policyholder},
\item
  the \textbf{direct insurer}, and
\item
  the \textbf{reinsurer}.
\end{itemize}

The way the loss is divided depends on the insurance and reinsurance
arrangements.

The main question is:

\begin{quote}
Given a gross claim, how much does each party pay?
\end{quote}
\end{frame}

\begin{frame}{The Basic Risk-Sharing Framework}
\protect\phantomsection\label{the-basic-risk-sharing-framework}
Let \(X\) denote the \textbf{gross claim amount}.

Define:

\begin{itemize}
\tightlist
\item
  \(V\): amount paid by the \textbf{policyholder};
\item
  \(Y\): amount ultimately paid by the \textbf{direct insurer};
\item
  \(Z\): amount paid by the \textbf{reinsurer}.
\end{itemize}

The payments must satisfy

\[
\boxed{X=V+Y+Z}.
\]

This is the basic accounting identity for risk sharing.
\end{frame}

\begin{frame}{The Three-Party Risk-Sharing Framework}
\protect\phantomsection\label{the-three-party-risk-sharing-framework}
\begin{figure}

\centering{

\includegraphics[width=0.85\linewidth,height=\textheight,keepaspectratio]{../figures/Basic_Risk_Sharing_Three-Party_Framework.png}

}

\caption{\label{fig-basic-risk-sharing}Basic Risk-Sharing Three-Party
Framework}

\end{figure}%
\end{frame}

\begin{frame}{No Deductible}
\protect\phantomsection\label{no-deductible}
Initially, suppose there is \textbf{no deductible}.

Then the policyholder does not contribute to the claim payment:

\[
V=0.
\]

Therefore,

\[
\boxed{X=Y+Z}.
\]

The gross claim is divided between:

\begin{itemize}
\tightlist
\item
  the amount retained by the insurer, \(Y\); and
\item
  the amount transferred to the reinsurer, \(Z\).
\end{itemize}
\end{frame}

\begin{frame}{Two Stages of Risk Sharing}
\protect\phantomsection\label{two-stages-of-risk-sharing}
Risk sharing can occur at two stages.

\textbf{Insurance}

The insurance contract determines how the loss is divided between:

\begin{itemize}
\tightlist
\item
  the policyholder, and
\item
  the direct insurer.
\end{itemize}

A deductible is one example.

\textbf{Reinsurance}

The direct insurer may transfer part of its risk to a reinsurer.

The reinsurer then pays part of the claim according to the reinsurance
contract.

Thus,

\[
\boxed{X=V+Y+Z}.
\]
\end{frame}

\begin{frame}{Why Does an Insurer Reinsure?}
\protect\phantomsection\label{why-does-an-insurer-reinsure}
An insurer may receive many claims of different sizes.

Most claims may be relatively small, but an individual claim can
occasionally be very large.

Without reinsurance:

\[
Y=X,
\qquad
Z=0.
\]

With reinsurance, part of the claim may be recovered from the reinsurer:

\[
Y<X,
\qquad
Z>0.
\]

Reinsurance therefore transfers part of the insurer's claim risk to
another insurer.
\end{frame}

\begin{frame}{Gross and Net Claims}
\protect\phantomsection\label{gross-and-net-claims}
The \textbf{gross claim} \(X\) is the loss before considering
reinsurance.

The \textbf{net claim} \(Y\) is the amount ultimately paid by the direct
insurer.

The \textbf{reinsurance recovery} \(Z\) is the amount paid by the
reinsurer.

With no deductible,

\[
X=Y+Z.
\]

Taking expectations,

\[
E(X)=E(Y)+E(Z).
\]

The same decomposition applies to other quantities, such as claim
variability.
\end{frame}

\begin{frame}{Reinsurance and Claim Variability}
\protect\phantomsection\label{reinsurance-and-claim-variability}
The gross claim distribution describes the losses generated by the
insurance portfolio.

The net claim distribution describes the amounts that remain with the
insurer after reinsurance.

For example,

\[
\operatorname{Var}(X)
\]

measures the variability of gross claims, whereas

\[
\operatorname{Var}(Y)
\]

measures the variability of the insurer's claim payments.

A purpose of reinsurance is to reduce the insurer's exposure to large
claims and hence reduce the variability of its claim payments.
\end{frame}

\begin{frame}{Two Broad Forms of Reinsurance}
\protect\phantomsection\label{two-broad-forms-of-reinsurance}
Reinsurance arrangements can be broadly classified as:

\textbf{Proportional reinsurance}

The reinsurer takes a fixed proportion of each claim.

\textbf{Non-proportional reinsurance}

The amount paid by the reinsurer depends on the size of the claim.

A common example is \textbf{excess-of-loss reinsurance}.

The key distinction is therefore:

\[
\text{fixed proportion}
\quad\text{versus}\quad
\text{claim size}.
\]
\end{frame}

\begin{frame}{Proportional versus Non-Proportional Reinsurance}
\protect\phantomsection\label{proportional-versus-non-proportional-reinsurance}
\begin{figure}

\centering{

\includegraphics[width=0.85\linewidth,height=\textheight,keepaspectratio]{../figures/Proportional_vs_Non-Proportional_Reinsurance.png}

}

\caption{\label{fig-reinsurance-comp}Proportional vs.~Non-Proportional
Reinsurance}

\end{figure}%
\end{frame}

\begin{frame}{Proportional Reinsurance}
\protect\phantomsection\label{proportional-reinsurance}
Let \(\alpha\) denote the \textbf{retained proportion}.

The insurer retains \(\alpha\) of every claim, while the reinsurer pays
\(1-\alpha\).

Thus,

\[
\boxed{Y=\alpha X} \quad \text{ and } \quad \boxed{Z=(1-\alpha)X}.
\] Hence,

\[
Y+Z=X.
\]

The insurer retains the same proportion regardless of the size of the
claim.
\end{frame}

\begin{frame}{Proportional Reinsurance: Example}
\protect\phantomsection\label{proportional-reinsurance-example}
Suppose the insurer retains \(75\%\) of every claim:

\[
\alpha=0.75.
\]

For a claim of \(100{,}000\),

\[
Y=0.75(100{,}000)=75{,}000,  \quad \text{ and } \quad    Z=0.25(100{,}000)=25{,}000.
\]

Thus,

\[
X=Y+Z.
\]

The same proportions apply to every claim.
\end{frame}

\begin{frame}{Excess-of-Loss Reinsurance}
\protect\phantomsection\label{excess-of-loss-reinsurance}
Suppose the insurer retains claims up to an amount \(M\).

The insurer pays:

\[
\boxed{Y=\min(X,M)}
\]

and the reinsurer pays:

\[
\boxed{Z=(X-M)_+},\quad \text{where } (x)_+=\max(x,0).
\] Again, \(X=Y+Z.\)

Unlike proportional reinsurance, the insurer's retained proportion is
not constant.
\end{frame}

\begin{frame}{Excess-of-Loss: How Does It Work?}
\protect\phantomsection\label{excess-of-loss-how-does-it-work}
If \(X\leq M\),

\[
Y=X,
\qquad
Z=0.
\]

The insurer pays the claim in full.

If \(X>M\),

\[
Y=M,
\qquad
Z=X-M.
\]

The insurer pays the retention and the reinsurer pays the excess.

Thus, large claims generate larger reinsurance recoveries.
\end{frame}

\begin{frame}{Excess-of-Loss: Example}
\protect\phantomsection\label{excess-of-loss-example}
Suppose the retention is

\[
M=50{,}000.
\]

{\def\LTcaptype{none} % do not increment counter
\begin{longtable}[]{@{}rrr@{}}
\toprule\noalign{}
Gross claim \(X\) & Insurer \(Y\) & Reinsurer \(Z\) \\
\midrule\noalign{}
\endhead
\(30{,}000\) & \(30{,}000\) & \(0\) \\
\(50{,}000\) & \(50{,}000\) & \(0\) \\
\(80{,}000\) & \(50{,}000\) & \(30{,}000\) \\
\(500{,}000\) & \(50{,}000\) & \(450{,}000\) \\
\bottomrule\noalign{}
\end{longtable}
}

For large claims, the insurer's payment is capped at \(M\).
\end{frame}

\begin{frame}{Comparing the Payment Functions}
\protect\phantomsection\label{comparing-the-payment-functions}
The two arrangements can be described directly through the payment
functions.

\textbf{Proportional}

\[
Y=\alpha X,
\qquad
Z=(1-\alpha)X.
\]

\textbf{Excess-of-loss}

\[
Y=\min(X,M),
\qquad
Z=(X-M)_+.
\]

The important question is always:

\begin{quote}
Given \(X\), how much does the insurer pay and how much does the
reinsurer pay?
\end{quote}
\end{frame}

\begin{frame}{A Unified View of Risk Sharing}
\protect\phantomsection\label{a-unified-view-of-risk-sharing}
The different arrangements can all be described through payment
functions.

For proportional reinsurance,

\[
\begin{aligned}
V&=0,\\
Y&=\alpha X,\\
Z&=(1-\alpha)X.
\end{aligned}
\]

For excess-of-loss reinsurance,

\[
\begin{aligned}
V&=0,\\
Y&=\min(X,M),\\
Z&=(X-M)_+.
\end{aligned}
\]

In both cases, \(\boxed{X=V+Y+Z}.\)
\end{frame}

\begin{frame}{From Claims to Payment Functions}
\protect\phantomsection\label{from-claims-to-payment-functions}
The central idea for the remainder of the chapter is:

\[
\boxed{
\text{Gross claim}
\longrightarrow
\text{risk-sharing arrangement}
\longrightarrow
\text{payment functions}
}
\]

Once the payment functions are known, we can analyse quantities such as

\[
E(Y),\qquad E(Z),\qquad
\operatorname{Var}(Y),
\]

using the distribution of the gross claim \(X\).
\end{frame}

\section{Reinsurance}\label{reinsurance}

\begin{frame}{Types of Reinsurance}
\protect\phantomsection\label{types-of-reinsurance}
Reinsurance is an arrangement under which an insurer transfers part of
the risk arising from its insurance portfolio to another insurer, known
as the \textbf{reinsurer}.

The original insurer is the \textbf{direct insurer} or \textbf{cedant}.

With no deductible,

\[
X=Y+Z,
\]

where:

\begin{itemize}
\tightlist
\item
  \(X\) is the gross claim;
\item
  \(Y\) is the amount ultimately paid by the direct insurer;
\item
  \(Z\) is the amount paid by the reinsurer.
\end{itemize}

The reinsurance contract determines how \(X\) is divided between \(Y\)
and \(Z\).
\end{frame}

\begin{frame}{Proportional Reinsurance}
\protect\phantomsection\label{proportional-reinsurance-1}
Under proportional reinsurance, the direct insurer and reinsurer share
each claim according to an agreed proportion.

Let \(\alpha\) denote the \textbf{retained proportion}.

Then

\[
\boxed{Y=\alpha X}
\]

and

\[
\boxed{Z=(1-\alpha)X}.
\]

Therefore,

\[
X=Y+Z.
\]

The key feature is

\[
\boxed{\frac{Y}{X}=\alpha},
\qquad X>0.
\]

The insurer retains the same proportion of every claim.
\end{frame}

\begin{frame}{Proportional Reinsurance: Example}
\protect\phantomsection\label{proportional-reinsurance-example-1}
Suppose the insurer retains \(75\%\) of each claim:

\[
\alpha=0.75.
\]

For a claim of \(100{,}000\),

\[
Y=0.75(100{,}000)=75{,}000,
\]

and

\[
Z=0.25(100{,}000)=25{,}000.
\]

Thus,

\[
X=75{,}000+25{,}000.
\]

The same proportions apply regardless of the claim size.
\end{frame}

\begin{frame}{Quota Share Reinsurance}
\protect\phantomsection\label{quota-share-reinsurance}
A common form of proportional reinsurance is \textbf{quota share
reinsurance}.

Under a quota share arrangement, the same retained proportion applies to
all risks in the portfolio.

For example, a \(75\%\) retained quota share means:

\[
Y=0.75X,
\qquad
Z=0.25X.
\]

Thus, the direct insurer retains \(75\%\) of every claim and the
reinsurer pays \(25\%\).

In this course, we focus on quota share reinsurance.
\end{frame}

\begin{frame}{Surplus Reinsurance}
\protect\phantomsection\label{surplus-reinsurance}
Another form of proportional reinsurance is \textbf{surplus
reinsurance}.

Under surplus reinsurance, the proportion ceded to the reinsurer can
vary from one risk to another.

Thus, unlike quota share reinsurance,

\[
\frac{Y}{X}
\]

does not have to be the same for every risk.

The common feature is that, for a given risk, the claim is divided
between the insurer and reinsurer in a fixed proportion.
\end{frame}

\begin{frame}{Excess-of-Loss Reinsurance}
\protect\phantomsection\label{excess-of-loss-reinsurance-1}
Under \textbf{excess-of-loss reinsurance}, the reinsurer does not pay a
fixed proportion of every claim.

Instead, the reinsurer pays the amount of a claim above a specified
\textbf{retention level}.

Let \(M\) denote the retention.

The insurer's payment is

\[
\boxed{Y=\min(X,M)}
\]

and the reinsurer's recovery is

\[
\boxed{Z=(X-M)_+},
\]

where

\[
(x)_+=\max(x,0).
\]

Again,

\[
X=Y+Z.
\]
\end{frame}

\begin{frame}{Excess-of-Loss: Two Cases}
\protect\phantomsection\label{excess-of-loss-two-cases}
If \(X\leq M\),

\[
Y=X,
\qquad
Z=0.
\]

The insurer pays the claim in full.

If \(X>M\),

\[
Y=M,
\qquad
Z=X-M.
\]

The insurer pays the retention, and the reinsurer pays the excess.

Thus, the insurer's payment is capped at \(M\).
\end{frame}

\begin{frame}{Excess-of-Loss: Example}
\protect\phantomsection\label{excess-of-loss-example-1}
Suppose the retention is

\[
M=50{,}000.
\]

{\def\LTcaptype{none} % do not increment counter
\begin{longtable}[]{@{}rrr@{}}
\toprule\noalign{}
Gross claim \(X\) & Insurer \(Y\) & Reinsurer \(Z\) \\
\midrule\noalign{}
\endhead
\(30{,}000\) & \(30{,}000\) & \(0\) \\
\(50{,}000\) & \(50{,}000\) & \(0\) \\
\(80{,}000\) & \(50{,}000\) & \(30{,}000\) \\
\(500{,}000\) & \(50{,}000\) & \(450{,}000\) \\
\bottomrule\noalign{}
\end{longtable}
}

For claims above \(50{,}000\), the insurer's payment does not increase.
\end{frame}

\begin{frame}{Reinsurance Layers}
\protect\phantomsection\label{reinsurance-layers}
An excess-of-loss contract may also specify a maximum amount that the
reinsurer will pay.

Suppose:

\begin{itemize}
\tightlist
\item
  \(M\) is the insurer's retention;
\item
  \(L\) is the maximum amount paid by the reinsurer.
\end{itemize}

The reinsurer's payment is

\[
\boxed{
Z=\min\{(X-M)_+,L\}.
}
\]

The cover therefore applies to the layer from

\[
M
\quad\text{to}\quad
M+L.
\]
\end{frame}

\begin{frame}{A Layer of Reinsurance}
\protect\phantomsection\label{a-layer-of-reinsurance}
Consider a layer of \(30{,}000\) excess of \(20{,}000\).

Then

\[
M=20{,}000,
\qquad
L=30{,}000,
\]

and

\[
Z=\min\{(X-20{,}000)_+,30{,}000\}.
\]

The reinsurer covers the part of the claim between

\[
20{,}000
\quad\text{and}\quad
50{,}000.
\]
\end{frame}

\begin{frame}{Reinsurance Layer: Example}
\protect\phantomsection\label{reinsurance-layer-example}
{\def\LTcaptype{none} % do not increment counter
\begin{longtable}[]{@{}rrr@{}}
\toprule\noalign{}
Gross claim \(X\) & Insurer \(Y\) & Reinsurer \(Z\) \\
\midrule\noalign{}
\endhead
\(15{,}000\) & \(15{,}000\) & \(0\) \\
\(30{,}000\) & \(20{,}000\) & \(10{,}000\) \\
\(50{,}000\) & \(20{,}000\) & \(30{,}000\) \\
\(80{,}000\) & \(50{,}000\) & \(30{,}000\) \\
\bottomrule\noalign{}
\end{longtable}
}

The reinsurer pays:

\begin{itemize}
\tightlist
\item
  nothing below the retention;
\item
  the excess above the retention within the layer;
\item
  the maximum \(L\) once the claim exceeds the layer.
\end{itemize}
\end{frame}

\begin{frame}{Retention, Limit, and Layer}
\protect\phantomsection\label{retention-limit-and-layer}
For a reinsurance layer:

\begin{itemize}
\tightlist
\item
  \textbf{Retention \(M\)}: the amount retained by the insurer before
  reinsurance begins.
\item
  \textbf{Limit \(L\)}: the maximum amount paid by the reinsurer.
\item
  \textbf{Layer}: the range of claim amounts covered by the reinsurer,
  from \(M\) to \(M+L\).
\end{itemize}

Thus,

\[
\boxed{
Z=\min\{(X-M)_+,L\}.
}
\]

If there is no upper limit, the cover is \textbf{unlimited}:

\[
Z=(X-M)_+.
\]
\end{frame}

\begin{frame}{Proportional versus Excess-of-Loss}
\protect\phantomsection\label{proportional-versus-excess-of-loss}
{\def\LTcaptype{none} % do not increment counter
\begin{longtable}[]{@{}
  >{\raggedright\arraybackslash}p{(\linewidth - 4\tabcolsep) * \real{0.3333}}
  >{\raggedright\arraybackslash}p{(\linewidth - 4\tabcolsep) * \real{0.3333}}
  >{\raggedright\arraybackslash}p{(\linewidth - 4\tabcolsep) * \real{0.3333}}@{}}
\toprule\noalign{}
\begin{minipage}[b]{\linewidth}\raggedright
\end{minipage} & \begin{minipage}[b]{\linewidth}\raggedright
Proportional
\end{minipage} & \begin{minipage}[b]{\linewidth}\raggedright
Excess-of-loss
\end{minipage} \\
\midrule\noalign{}
\endhead
Basis of payment & Fixed proportion of claim & Amount above retention \\
Insurer payment & \(Y=\alpha X\) & \(Y=\min(X,M)\) \\
Reinsurer payment & \(Z=(1-\alpha)X\) & \(Z=(X-M)_+\) \\
Small claims & Shared & Usually retained \\
Large claims & Shared in the same proportion & Substantial part
transferred \\
\bottomrule\noalign{}
\end{longtable}
}

The main difference is how the reinsurer's payment responds to the size
of the claim.
\end{frame}

\begin{frame}{What Is the Main Difference?}
\protect\phantomsection\label{what-is-the-main-difference}
Under proportional reinsurance,

\[
Z=(1-\alpha)X.
\]

The reinsurer shares \textbf{every claim} in a fixed proportion.

Under excess-of-loss reinsurance,

\[
Z=(X-M)_+.
\]

The reinsurer primarily protects the insurer against \textbf{larger
claims}.

This distinction affects both:
\(E(Y) \quad\text{and}\quad \operatorname{Var}(Y).\)

The payment functions allow these quantities to be analysed
mathematically.
\end{frame}

\begin{frame}
\end{frame}

\section{Deductible and Policy}\label{deductible-and-policy}

\begin{frame}{From Reinsurance to Deductibles}
\protect\phantomsection\label{from-reinsurance-to-deductibles}
Reinsurance transfers part of the insurer's risk to the
\textbf{reinsurer}.

A deductible transfers part of the risk to the \textbf{policyholder}.

Both can be described using the same risk-sharing framework.

Without reinsurance,

\[
X=V+Y.
\]

With reinsurance,

\[
\boxed{X=V+Y+Z}.
\]

We now focus on the relationship between: \(X,\qquad V,\qquad Y.\)
\end{frame}

\begin{frame}{Deductibles}
\protect\phantomsection\label{deductibles}
A \textbf{deductible} is an amount of a claim borne by the policyholder
before the insurer makes a payment.

Suppose the deductible is \(d\).

The policyholder pays

\[
\boxed{V=\min(X,d)}
\]

and the insurer pays

\[
\boxed{Y=(X-d)_+}.
\]

Therefore,

\[
\boxed{X=V+Y}.
\]
\end{frame}

\begin{frame}{Ordinary Deductible: Two Cases}
\protect\phantomsection\label{ordinary-deductible-two-cases}
If \(X\leq d\),

\[
V=X,
\qquad
Y=0.
\]

The policyholder pays the claim in full.

If \(X>d\),

\[
V=d,
\qquad
Y=X-d.
\]

The policyholder pays the deductible and the insurer pays the remainder.
\end{frame}

\begin{frame}{Ordinary Deductible: Example}
\protect\phantomsection\label{ordinary-deductible-example}
Suppose

\[
d=10{,}000.
\]

{\def\LTcaptype{none} % do not increment counter
\begin{longtable}[]{@{}rrr@{}}
\toprule\noalign{}
Gross claim \(X\) & Policyholder \(V\) & Insurer \(Y\) \\
\midrule\noalign{}
\endhead
\(5{,}000\) & \(5{,}000\) & \(0\) \\
\(10{,}000\) & \(10{,}000\) & \(0\) \\
\(25{,}000\) & \(10{,}000\) & \(15{,}000\) \\
\(80{,}000\) & \(10{,}000\) & \(70{,}000\) \\
\bottomrule\noalign{}
\end{longtable}
}

The deductible eliminates insurer payments for claims below \(d\).

For larger claims, the policyholder's payment is capped at \(d\).
\end{frame}

\begin{frame}{Policy Limits}
\protect\phantomsection\label{policy-limits}
A \textbf{policy limit} is the maximum amount that the insurer will pay
for a claim.

Suppose the policy limit is \(u\) and there is no deductible.

The insurer pays

\[
\boxed{Y=\min(X,u)}.
\]

The policyholder pays the amount above the limit:

\[
\boxed{V=(X-u)_+}.
\]

Hence,

\[
X=V+Y.
\]
\end{frame}

\begin{frame}{Policy Limit: Example}
\protect\phantomsection\label{policy-limit-example}
Suppose

\[
u=100{,}000.
\]

{\def\LTcaptype{none} % do not increment counter
\begin{longtable}[]{@{}rrr@{}}
\toprule\noalign{}
Gross claim \(X\) & Policyholder \(V\) & Insurer \(Y\) \\
\midrule\noalign{}
\endhead
\(50{,}000\) & \(0\) & \(50{,}000\) \\
\(100{,}000\) & \(0\) & \(100{,}000\) \\
\(150{,}000\) & \(50{,}000\) & \(100{,}000\) \\
\(300{,}000\) & \(200{,}000\) & \(100{,}000\) \\
\bottomrule\noalign{}
\end{longtable}
}

The insurer's payment is capped at \(u\).
\end{frame}

\begin{frame}{Deductible and Policy Limit}
\protect\phantomsection\label{deductible-and-policy-limit}
Suppose a policy has:

\begin{itemize}
\tightlist
\item
  deductible \(d\);
\item
  policy limit \(u\), where \(u>d\).
\end{itemize}

The insurer pays the part of the claim between \(d\) and \(u\):

\[
\boxed{
Y=\min\{(X-d)_+,u-d\}.
}
\]

Equivalently,

\[
Y=
\begin{cases}
0, & X\leq d,\\
X-d, & d<X\leq u,\\
u-d, & X>u.
\end{cases}
\]

The maximum insurer payment is

\[
\boxed{u-d}.
\]
\end{frame}

\begin{frame}{Deductible and Policy Limit: Example}
\protect\phantomsection\label{deductible-and-policy-limit-example}
Suppose

\[
d=10{,}000,
\qquad
u=100{,}000.
\]

{\def\LTcaptype{none} % do not increment counter
\begin{longtable}[]{@{}rrr@{}}
\toprule\noalign{}
Gross claim \(X\) & Policyholder \(V\) & Insurer \(Y\) \\
\midrule\noalign{}
\endhead
\(5{,}000\) & \(5{,}000\) & \(0\) \\
\(25{,}000\) & \(10{,}000\) & \(15{,}000\) \\
\(100{,}000\) & \(10{,}000\) & \(90{,}000\) \\
\(150{,}000\) & \(60{,}000\) & \(90{,}000\) \\
\bottomrule\noalign{}
\end{longtable}
}

The deductible determines where the insurer's payment begins.

The policy limit determines where the insurer's payment stops
increasing.
\end{frame}

\begin{frame}{A Unified Risk-Sharing View}
\protect\phantomsection\label{a-unified-risk-sharing-view}
The arrangements considered so far can all be described through payment
functions.

\textbf{Deductible}

\[
Y=(X-d)_+.
\]

\textbf{Policy limit}

\[
Y=\min(X,u).
\]

\textbf{Deductible and policy limit}

\[
Y=\min\{(X-d)_+,u-d\}.
\]

\textbf{Excess-of-loss reinsurance}

\[
Z=(X-M)_+.
\]

The same mathematical structure appears repeatedly.
\end{frame}

\begin{frame}{Policyholder, Insurer, and Reinsurer}
\protect\phantomsection\label{policyholder-insurer-and-reinsurer}
The full risk-sharing relationship is

\[
\boxed{X=V+Y+Z}.
\]

Each contract determines a payment function:

\[
X
\longrightarrow
\begin{cases}
V & \text{policyholder},\\
Y & \text{direct insurer},\\
Z & \text{reinsurer}.
\end{cases}
\]

This gives a common framework for analysing different insurance and
reinsurance arrangements.
\end{frame}

\begin{frame}{Payment Functions and Distributions}
\protect\phantomsection\label{payment-functions-and-distributions}
The payment functions determine the distributions of the amounts
actually paid.

For example, under excess-of-loss reinsurance,

\[
Y=\min(X,M),
\qquad
Z=(X-M)_+.
\]

Even if \(X\) has a continuous distribution, \(Y\) and \(Z\) need not be
purely continuous.

Why?

Because the payment functions can map a range of values of \(X\) to the
same payment amount.

This produces \textbf{point masses} in the payment distributions.
\end{frame}

\begin{frame}{Insurer's Payment under Excess-of-Loss}
\protect\phantomsection\label{insurers-payment-under-excess-of-loss}
Suppose \(X\) is continuous with distribution function \(F_X\).

Under unlimited excess-of-loss reinsurance,

\[
Y=\min(X,M).
\]

The insurer pays the full claim when \(X\leq M\) and exactly \(M\) when
\(X>M\).

For \(y<M\),

\[
F_Y(y)=F_X(y).
\]

For \(y\geq M\),

\[
F_Y(y)=1.
\]

Thus,

\[
\boxed{
F_Y(y)=
\begin{cases}
F_X(y), & y<M,\\
1, & y\geq M.
\end{cases}
}
\]
\end{frame}

\begin{frame}{A Point Mass at the Retention}
\protect\phantomsection\label{a-point-mass-at-the-retention}
All claims satisfying \(X>M\) result in

\[
Y=M.
\]

Therefore,

\[
\boxed{
\Pr(Y=M)=\Pr(X>M)=1-F_X(M).
}
\]

Thus, even if \(X\) is continuous, \(Y\) has a positive probability mass
at \(M\).

The distribution of \(Y\) consists of:

\[
\text{continuous component on }(0,M)
\]

and

\[
\text{point mass at }M.
\]
\end{frame}

\begin{frame}
\begin{tcolorbox}[enhanced jigsaw, arc=.35mm, bottomrule=.15mm, bottomtitle=1mm, breakable, colback=white, colbacktitle=quarto-callout-note-color!10!white, colframe=quarto-callout-note-color-frame, coltitle=black, left=2mm, leftrule=.75mm, opacityback=0, opacitybacktitle=0.6, rightrule=.15mm, title=\textcolor{quarto-callout-note-color}{\faInfo}\hspace{0.5em}{Key idea}, titlerule=0mm, toprule=.15mm, toptitle=1mm]

A payment limit can create a \textbf{point mass}.

Under excess-of-loss reinsurance,

\[
Y=\min(X,M)
\]

maps all claims exceeding \(M\) to the single payment \(M\).

Hence,

\[
\Pr(Y=M)>0.
\]

\end{tcolorbox}
\end{frame}

\begin{frame}{Example: Distribution of the Insurer's Payment}
\protect\phantomsection\label{example-distribution-of-the-insurers-payment}
Let

\[
X\sim\operatorname{Exponential}(\lambda)
\]

and suppose the insurer has unlimited excess-of-loss reinsurance with
retention \(M\).

The insurer's payment is

\[
Y=\min(X,M).
\]

Find:

\begin{enumerate}
\tightlist
\item
  the distribution function \(F_Y(y)\);
\item
  the continuous component of the distribution;
\item
  the probability mass at \(M\);
\item
  \(\Pr(Y=M)\).
\end{enumerate}

This example illustrates how a risk-sharing arrangement changes the
distribution of a claim payment.
\end{frame}

\begin{frame}{Reinsurer's Payment under Excess-of-Loss}
\protect\phantomsection\label{reinsurers-payment-under-excess-of-loss}
The reinsurer pays

\[
Z=(X-M)_+.
\]

Thus,

\begin{itemize}
\tightlist
\item
  if \(X\leq M\), then \(Z=0\);
\item
  if \(X>M\), then \(Z=X-M\).
\end{itemize}

The reinsurer therefore has a positive probability of making \textbf{no
payment}.

At zero,

\[
\boxed{
\Pr(Z=0)=\Pr(X\leq M)=F_X(M).
}
\]
\end{frame}

\begin{frame}{Distribution of the Reinsurer's Payment}
\protect\phantomsection\label{distribution-of-the-reinsurers-payment}
For \(z>0\),

\[
\begin{aligned}
F_Z(z)
&=\Pr(Z\leq z)\\
&=\Pr((X-M)_+\leq z)\\
&=\Pr(X\leq M+z)\\
&=F_X(M+z).
\end{aligned}
\]

Hence,

\[
\boxed{
F_Z(z)=
\begin{cases}
0, & z<0,\\
F_X(M), & z=0,\\
F_X(M+z), & z>0.
\end{cases}
}
\]

The distribution has a point mass at zero.
\end{frame}

\begin{frame}{Two Different Point Masses}
\protect\phantomsection\label{two-different-point-masses}
Under excess-of-loss reinsurance,

\[
Y=\min(X,M)
\]

has a point mass at the \textbf{retention \(M\)}:

\[
\Pr(Y=M)=1-F_X(M).
\]

Meanwhile,

\[
Z=(X-M)_+
\]

has a point mass at \textbf{zero}:

\[
\Pr(Z=0)=F_X(M).
\]

The two point masses arise from the same reinsurance arrangement.
\end{frame}

\begin{frame}{Mixed Distributions}
\protect\phantomsection\label{mixed-distributions}
A random variable can have both:

\begin{itemize}
\tightlist
\item
  a continuous component; and
\item
  one or more point masses.
\end{itemize}

Such a random variable is called a \textbf{mixed distribution}.

For example, \(Y=\min(X,M)\) has:

\[
\text{continuous component on }(0,M)
\]

and

\[
\Pr(Y=M)=1-F_X(M).
\]

Similarly, \(Z=(X-M)_+\) has:

\[
\Pr(Z=0)=F_X(M)
\]

and a continuous component for \(z>0\).
\end{frame}

\begin{frame}{Reinsurance Claims}
\protect\phantomsection\label{reinsurance-claims}
The reinsurer's payment

\[
Z=(X-M)_+
\]

includes all claims, including claims for which the reinsurer pays zero.

In practice, we may instead consider only claims that actually involve
the reinsurance contract:

\[
X>M.
\]

These are \textbf{reinsurance claims}.

It is therefore useful to distinguish between:

\[
Z
\quad\text{and}\quad
W=Z\mid Z>0.
\]
\end{frame}

\begin{frame}{Payment Conditional on a Reinsurance Claim}
\protect\phantomsection\label{payment-conditional-on-a-reinsurance-claim}
Since

\[
Z>0
\iff
X>M,
\]

we have

\[
W=X-M\mid X>M.
\]

For \(w>0\),

\[
\begin{aligned}
F_W(w)
&=\Pr(X-M\leq w\mid X>M)\\
&=\frac{\Pr(M<X\leq M+w)}
{\Pr(X>M)}\\
&=\frac{F_X(M+w)-F_X(M)}
{1-F_X(M)}.
\end{aligned}
\]

Thus, \(W\) describes the size of a \textbf{positive reinsurance
payment}.
\end{frame}

\begin{frame}{Three Claim Amounts to Distinguish}
\protect\phantomsection\label{three-claim-amounts-to-distinguish}
It is useful to keep the following random variables separate:

\[
\begin{aligned}
X &: \text{gross claim amount},\\
Z &: \text{reinsurer's payment on a randomly selected claim},\\
W &: \text{reinsurer's payment given that the claim involves reinsurance}.
\end{aligned}
\]

The distinction between \(Z\) and \(W\) is important.

\(Z\) has a point mass at zero.

\(W\) describes only positive reinsurance payments.
\end{frame}

\begin{frame}{Expected Payments}
\protect\phantomsection\label{expected-payments}
The payment functions can be used to calculate expected payments.

For the insurer,

\[
Y=\min(X,M).
\]

Therefore,

\[
E[Y]
=
\int_0^M x f_X(x)\,dx
+
M\{1-F_X(M)\}.
\]

The second term accounts for the point mass at \(M\).

For the reinsurer,

\[
Z=(X-M)_+,
\]

so

\[
\boxed{
E[Z]
=
\int_M^\infty (x-M)f_X(x)\,dx.
}
\]
\end{frame}

\begin{frame}{Expected Payments: Risk-Sharing Identity}
\protect\phantomsection\label{expected-payments-risk-sharing-identity}
Since

\[
X=Y+Z,
\]

taking expectations gives

\[
\boxed{
E[X]=E[Y]+E[Z].
}
\]

Equivalently,

\[
E[Y]=E[X]-E[Z].
\]

Thus, the expected gross claim is divided between:

\begin{itemize}
\tightlist
\item
  the expected amount retained by the insurer; and
\item
  the expected amount transferred to the reinsurer.
\end{itemize}

Reinsurance changes \textbf{who bears the loss}.
\end{frame}

\begin{frame}{Exponential Claims}
\protect\phantomsection\label{exponential-claims}
Suppose

\[
X\sim\operatorname{Exponential}(\lambda).
\]

Then

\[
E[X]=\frac{1}{\lambda}.
\]

Under excess-of-loss reinsurance with retention \(M\),

\[
\Pr(X>M)=e^{-\lambda M}.
\]

Thus, the proportion of claims involving the reinsurer is

\[
\boxed{e^{-\lambda M}}.
\]
\end{frame}

\begin{frame}{Exponential Claims: Expected Payments}
\protect\phantomsection\label{exponential-claims-expected-payments}
For the reinsurer,

\[
E[Z]
=
\int_M^\infty (x-M)\lambda e^{-\lambda x}\,dx.
\]

This gives

\[
\boxed{
E[Z]=\frac{e^{-\lambda M}}{\lambda}.
}
\]

Since

\[
E[X]=\frac{1}{\lambda},
\]

we can write

\[
\boxed{
E[Z]=e^{-\lambda M}E[X].
}
\]

Therefore,

\[
\boxed{
E[Y]=\frac{1-e^{-\lambda M}}{\lambda}.
}
\]
\end{frame}

\begin{frame}{Effect of the Retention}
\protect\phantomsection\label{effect-of-the-retention}
For excess-of-loss reinsurance,

\[
Y=\min(X,M).
\]

As the retention \(M\) increases:

\begin{itemize}
\tightlist
\item
  fewer claims involve the reinsurer;
\item
  the insurer retains more of the expected claim;
\item
  the reinsurer pays less on average.
\end{itemize}

For exponential claims,

\[
\Pr(X>M)=e^{-\lambda M}
\]

and

\[
E[Z]=\frac{e^{-\lambda M}}{\lambda}.
\]

Thus, the retention directly controls the extent of risk transferred to
the reinsurer.
\end{frame}

\begin{frame}{Deductibles and Policy Limits}
\protect\phantomsection\label{deductibles-and-policy-limits}
Reinsurance transfers part of the insurer's risk to a reinsurer.

A \textbf{deductible} transfers part of the risk to the policyholder.

A \textbf{policy limit} restricts the maximum amount that the insurer
will pay.

Using the same notation,

\[
X=V+Y+Z.
\]

In the absence of reinsurance,

\[
\boxed{X=V+Y}.
\]

In this section, we focus on the relationship between:

\[
X,\qquad V,\qquad Y.
\]
\end{frame}

\begin{frame}{Ordinary Deductible}
\protect\phantomsection\label{ordinary-deductible}
Suppose a policy has an ordinary deductible of amount \(d\).

The policyholder pays the first \(d\) of the claim, while the insurer
pays the amount exceeding \(d\).

Therefore,

\[
\boxed{V=\min(X,d)}
\]

and

\[
\boxed{Y=(X-d)_+}.
\]

Since

\[
\min(X,d)+(X-d)_+=X,
\]

we have

\[
\boxed{X=V+Y}.
\]
\end{frame}

\begin{frame}{Ordinary Deductible: How Does It Work?}
\protect\phantomsection\label{ordinary-deductible-how-does-it-work}
There are two cases.

If \(X\leq d\),

\[
V=X,
\qquad
Y=0.
\]

The claim is entirely below the deductible.

If \(X>d\),

\[
V=d,
\qquad
Y=X-d.
\]

The policyholder pays the deductible and the insurer pays the remainder.
\end{frame}

\begin{frame}{Ordinary Deductible: Example}
\protect\phantomsection\label{ordinary-deductible-example-1}
Suppose

\[
d=10{,}000.
\]

{\def\LTcaptype{none} % do not increment counter
\begin{longtable}[]{@{}rrr@{}}
\toprule\noalign{}
Gross claim \(X\) & Policyholder \(V\) & Insurer \(Y\) \\
\midrule\noalign{}
\endhead
\(5{,}000\) & \(5{,}000\) & \(0\) \\
\(10{,}000\) & \(10{,}000\) & \(0\) \\
\(25{,}000\) & \(10{,}000\) & \(15{,}000\) \\
\(80{,}000\) & \(10{,}000\) & \(70{,}000\) \\
\bottomrule\noalign{}
\end{longtable}
}

The deductible has two effects:

\begin{itemize}
\tightlist
\item
  it completely eliminates insurer payments for claims below \(d\);
\item
  for larger claims, it reduces the insurer's payment by \(d\).
\end{itemize}
\end{frame}

\begin{frame}{Policy Limit}
\protect\phantomsection\label{policy-limit}
A \textbf{policy limit} is the maximum amount that the insurer will pay
for a claim.

Suppose there is a policy limit of \(u\) and no deductible.

The insurer pays the claim in full when

\[
X\leq u,
\]

but pays only \(u\) when

\[
X>u.
\]

Thus,

\[
\boxed{Y=\min(X,u)}.
\]

The amount above the policy limit is borne by the policyholder:

\[
\boxed{V=(X-u)_+}.
\]

Hence,

\[
X=V+Y.
\]
\end{frame}

\begin{frame}{Policy Limit: Example}
\protect\phantomsection\label{policy-limit-example-1}
Suppose

\[
u=100{,}000.
\]

{\def\LTcaptype{none} % do not increment counter
\begin{longtable}[]{@{}rrr@{}}
\toprule\noalign{}
Gross claim \(X\) & Policyholder \(V\) & Insurer \(Y\) \\
\midrule\noalign{}
\endhead
\(50{,}000\) & \(0\) & \(50{,}000\) \\
\(100{,}000\) & \(0\) & \(100{,}000\) \\
\(150{,}000\) & \(50{,}000\) & \(100{,}000\) \\
\(300{,}000\) & \(200{,}000\) & \(100{,}000\) \\
\bottomrule\noalign{}
\end{longtable}
}

The insurer's payment is capped at \(u\).

This protects the insurer against very large claims.
\end{frame}

\begin{frame}{Deductible and Policy Limit Together}
\protect\phantomsection\label{deductible-and-policy-limit-together}
Suppose a policy has:

\begin{itemize}
\tightlist
\item
  deductible \(d\);
\item
  policy limit \(u\), where \(u>d\).
\end{itemize}

The insurer pays the part of the claim between \(d\) and \(u\).

Therefore,

\[
\boxed{
Y=\min\{(X-d)_+,u-d\}.
}
\]

Equivalently,

\[
Y=\min(X,u)-\min(X,d).
\]

The policyholder pays the remainder:

\[
\boxed{V=X-Y}.
\]
\end{frame}

\begin{frame}{Deductible and Policy Limit: Payment Function}
\protect\phantomsection\label{deductible-and-policy-limit-payment-function}
The insurer's payment can be written as

\[
Y=
\begin{cases}
0, & X\leq d,\\
X-d, & d<X\leq u,\\
u-d, & X>u.
\end{cases}
\]

The payment therefore has three regions:

\[
\begin{array}{ccl}
X\leq d
&\longrightarrow&
Y=0,\\[4pt]
d<X\leq u
&\longrightarrow&
Y=X-d,\\[4pt]
X>u
&\longrightarrow&
Y=u-d.
\end{array}
\]

The insurer's maximum payment is

\[
\boxed{u-d}.
\]
\end{frame}

\begin{frame}{Deductible and Policy Limit: Example}
\protect\phantomsection\label{deductible-and-policy-limit-example-1}
Suppose

\[
d=10{,}000,
\qquad
u=100{,}000.
\]

{\def\LTcaptype{none} % do not increment counter
\begin{longtable}[]{@{}rrr@{}}
\toprule\noalign{}
Gross claim \(X\) & Policyholder \(V\) & Insurer \(Y\) \\
\midrule\noalign{}
\endhead
\(5{,}000\) & \(5{,}000\) & \(0\) \\
\(25{,}000\) & \(10{,}000\) & \(15{,}000\) \\
\(100{,}000\) & \(10{,}000\) & \(90{,}000\) \\
\(150{,}000\) & \(60{,}000\) & \(90{,}000\) \\
\bottomrule\noalign{}
\end{longtable}
}

For claims above the policy limit, the insurer continues to pay only
\(90{,}000\).

The policyholder bears both:

\begin{itemize}
\tightlist
\item
  the deductible; and
\item
  the amount above the policy limit.
\end{itemize}
\end{frame}

\begin{frame}{Interpreting the Deductible and Policy Limit}
\protect\phantomsection\label{interpreting-the-deductible-and-policy-limit}
The two parameters control different parts of the insurer's payment.

\textbf{Deductible \(d\)}

determines where the insurer's payment begins.

\textbf{Policy limit \(u\)}

determines where the insurer's payment stops increasing.

Thus,

\[
\boxed{
\text{Insurer's payment}
=
\text{claim amount within the layer }(d,u).
}
\]

The maximum insurer payment is

\[
u-d.
\]
\end{frame}

\begin{frame}{Risk-Sharing Interpretation}
\protect\phantomsection\label{risk-sharing-interpretation}
The deductible and policy limit can be viewed as a \textbf{layer of
insurance}.

For a claim \(X\):

\[
\boxed{
Y=\min\{(X-d)_+,u-d\}.
}
\]

The insurer covers the layer

\[
\boxed{(d,u]}.
\]

The policyholder bears:

\begin{itemize}
\tightlist
\item
  the layer below \(d\); and
\item
  the amount above \(u\).
\end{itemize}

This is the same layer structure that appears in excess-of-loss
reinsurance.
\end{frame}

\begin{frame}{Insurance and Reinsurance Layers}
\protect\phantomsection\label{insurance-and-reinsurance-layers}
The same mathematical form appears in both insurance and reinsurance.

For a policy with deductible \(d\) and limit \(u\),

\[
Y=\min\{(X-d)_+,u-d\}.
\]

For a reinsurance layer with retention \(M\) and limit \(L\),

\[
Z=\min\{(X-M)_+,L\}.
\]

The difference is mainly the \textbf{party receiving the payment}.

In the insurance contract:

\[
Y=\text{insurer's payment}.
\]

In the reinsurance contract:

\[
Z=\text{reinsurer's payment}.
\]
\end{frame}

\begin{frame}{Comparing Risk-Sharing Arrangements}
\protect\phantomsection\label{comparing-risk-sharing-arrangements}
The main payment functions introduced so far are:

{\def\LTcaptype{none} % do not increment counter
\begin{longtable}[]{@{}ll@{}}
\toprule\noalign{}
Arrangement & Payment \\
\midrule\noalign{}
\endhead
Proportional reinsurance & \(Y=\alpha X\) \\
Excess-of-loss reinsurance & \(Z=(X-M)_+\) \\
Ordinary deductible & \(Y=(X-d)_+\) \\
Policy limit & \(Y=\min(X,u)\) \\
Deductible and limit & \(Y=\min\{(X-d)_+,u-d\}\) \\
\bottomrule\noalign{}
\end{longtable}
}

Each arrangement can therefore be analysed by specifying a payment
function.
\end{frame}

\begin{frame}{The Common Mathematical Structure}
\protect\phantomsection\label{the-common-mathematical-structure}
The central quantity is always the gross claim \(X\).

A contract transforms \(X\) into the amount paid by a particular party.

For example,

\[
X
\longrightarrow
Y=\min\{(X-d)_+,u-d\}
\]

for an insurance layer, while

\[
X
\longrightarrow
Z=\min\{(X-M)_+,L\}
\]

for a reinsurance layer.

Once the payment function is known, we can study its distribution and
moments.
\end{frame}

\begin{frame}
\end{frame}

\section{Payment Distributions under Risk
Sharing}\label{payment-distributions-under-risk-sharing}

\begin{frame}{How Risk Sharing Changes the Claim Distribution}
\protect\phantomsection\label{how-risk-sharing-changes-the-claim-distribution}
Risk-sharing arrangements change the distribution of the amount paid by
each party.

Suppose \(X\) is the gross claim amount.

After applying a payment function,

\[
Y=g(X),
\]

the random variable \(Y\) may have a distribution that is quite
different from that of \(X\).

In particular, payment limits can create \textbf{point masses}.
\end{frame}

\begin{frame}{Insurer's Payment under Excess-of-Loss}
\protect\phantomsection\label{insurers-payment-under-excess-of-loss-1}
Consider unlimited excess-of-loss reinsurance with retention \(M\).

The insurer's payment is

\[
\boxed{Y=\min(X,M)}.
\]

There are two cases:

\[
Y=
\begin{cases}
X, & X\leq M,\\
M, & X>M.
\end{cases}
\]

Thus, every claim exceeding \(M\) produces exactly the same insurer
payment: \(Y=M.\)

This has an important consequence for the distribution of \(Y\).
\end{frame}

\begin{frame}{Why Is \(Y\) a Mixed Distribution?}
\protect\phantomsection\label{why-is-y-a-mixed-distribution}
Suppose \(X\) has a continuous distribution.

For \(X<M\), different values of \(X\) produce different values of
\(Y\).

However, all values

\[
X>M
\]

are mapped to the single value \(Y = M.\)

Therefore, \(\quad \quad \Pr(Y=M)=\Pr(X>M).\)

If \(\Pr(X>M)>0\), then \(Y\) has a positive probability mass at \(M\).

Hence, \(Y\) is a \textbf{mixed distribution}:

\[
\boxed{
\text{continuous part}
+
\text{point mass at }M.
}
\]
\end{frame}

\begin{frame}{Distribution Function of \(Y\)}
\protect\phantomsection\label{distribution-function-of-y}
For

\[
Y=\min(X,M),
\]

the distribution function is

\[
F_Y(y)=\Pr(Y\leq y).
\]

We consider two cases according to the value of \(y\) relative to \(M\).
\end{frame}

\begin{frame}{Distribution Function of \(Y\)}
\protect\phantomsection\label{distribution-function-of-y-1}
For \(y<M\),

\[
Y\leq y
\iff
X\leq y.
\]

Therefore,

\[
\boxed{
F_Y(y)=F_X(y),
\qquad y<M.
}
\]
\end{frame}

\begin{frame}{Distribution Function of \(Y\)}
\protect\phantomsection\label{distribution-function-of-y-2}
For \(y\geq M\),

\[
Y=\min(X,M)\leq M\leq y,
\]

so

\[
F_Y(y)=1.
\]

Combining the two cases,

\[
\boxed{
F_Y(y)=
\begin{cases}
F_X(y), & y<M,\\
1, & y\geq M.
\end{cases}
}
\]
\end{frame}

\begin{frame}{Point Mass at the Retention}
\protect\phantomsection\label{point-mass-at-the-retention}
The insurer's payment under excess-of-loss reinsurance is

\[
Y=\min(X,M).
\]

All claims satisfying

\[
X\geq M
\]

produce the payment

\[
Y=M.
\]

Therefore,

\[
\Pr(Y=M)=\Pr(X\geq M).
\]
\end{frame}

\begin{frame}{Point Mass at the Retention}
\protect\phantomsection\label{point-mass-at-the-retention-1}
If \(X\) is continuous, then

\[
\Pr(X=M)=0.
\]

Hence,

\[
\boxed{
\Pr(Y=M)
=
\Pr(X>M)
=
1-F_X(M).
}
\]

Thus, the payment limit creates a \textbf{point mass at the retention
\(M\)}.
\end{frame}

\begin{frame}
\begin{tcolorbox}[enhanced jigsaw, arc=.35mm, bottomrule=.15mm, bottomtitle=1mm, breakable, colback=white, colbacktitle=quarto-callout-note-color!10!white, colframe=quarto-callout-note-color-frame, coltitle=black, left=2mm, leftrule=.75mm, opacityback=0, opacitybacktitle=0.6, rightrule=.15mm, title=\textcolor{quarto-callout-note-color}{\faInfo}\hspace{0.5em}{Key idea}, titlerule=0mm, toprule=.15mm, toptitle=1mm]

A continuous gross claim does not necessarily produce a continuous
payment distribution.

If

\[
Y=\min(X,M),
\]

all claims above \(M\) produce the same payment \(Y=M\).

Therefore,

\[
\Pr(Y=M)>0
\]

whenever \(\Pr(X>M)>0\).

\end{tcolorbox}
\end{frame}

\begin{frame}{Example: Distribution of the Insurer's Payment}
\protect\phantomsection\label{example-distribution-of-the-insurers-payment-1}
Let

\[
X\sim\operatorname{Exponential}(\lambda),
\]

and suppose the insurer has excess-of-loss reinsurance with retention
\(M\).

The insurer's payment is

\[
Y=\min(X,M).
\]

Find:

\begin{enumerate}
\tightlist
\item
  the distribution function \(F_Y(y)\);
\item
  the continuous part of the distribution;
\item
  the probability mass at \(M\).
\end{enumerate}
\end{frame}

\begin{frame}{Distribution of the Insurer's Payment}
\protect\phantomsection\label{distribution-of-the-insurers-payment}
For \(y<M\),

\[
\begin{aligned}
F_Y(y)
&=\Pr(Y\leq y)\\
&=\Pr(X\leq y).
\end{aligned}
\]

Hence,

\[
F_Y(y)=F_X(y).
\]

For \(y\geq M\),

\[
F_Y(y)=1.
\]
\end{frame}

\begin{frame}{Distribution of the Insurer's Payment}
\protect\phantomsection\label{distribution-of-the-insurers-payment-1}
Therefore,

\[
\boxed{
F_Y(y)=
\begin{cases}
F_X(y), & y<M,\\
1, & y\geq M.
\end{cases}
}
\]

For the exponential distribution,

\[
F_X(y)=1-e^{-\lambda y},
\qquad y\geq0.
\]

Thus,

\[
F_Y(y)=
\begin{cases}
1-e^{-\lambda y}, & 0\leq y<M,\\
1, & y\geq M.
\end{cases}
\]
\end{frame}

\begin{frame}{The Point Mass at \(M\)}
\protect\phantomsection\label{the-point-mass-at-m}
All claims satisfying

\[
X>M
\]

produce the same insurer payment:

\[
Y=M.
\]

Therefore,

\[
\begin{aligned}
\Pr(Y=M)
&=\Pr(X>M) =1-F_X(M).
\end{aligned}
\]

For exponential claims,

\[
\boxed{
\Pr(Y=M)=e^{-\lambda M}.
}
\]
\end{frame}

\begin{frame}{Distribution of the Insurer's Payment}
\protect\phantomsection\label{distribution-of-the-insurers-payment-2}
The distribution of \(Y\) therefore has two components:

\[
\boxed{
\begin{aligned}
&\text{continuous part on }0\leq y<M,\\
&\text{point mass at }M.
\end{aligned}
}
\]

The point mass is

\[
\Pr(Y=M)=1-F_X(M).
\]

This is an example of a \textbf{mixed distribution}.
\end{frame}

\begin{frame}{Reinsurer's Payment}
\protect\phantomsection\label{reinsurers-payment}
Under the same excess-of-loss arrangement,

\[
\boxed{
Z=(X-M)_+.
}
\]

Equivalently,

\[
Z=
\begin{cases}
0, & X\leq M,\\
X-M, & X>M.
\end{cases}
\]

The reinsurer pays nothing when the claim does not exceed the retention.
\end{frame}

\begin{frame}{Point Mass at Zero}
\protect\phantomsection\label{point-mass-at-zero}
Since

\[
Z=0
\iff
X\leq M,
\]

we have

\[
\boxed{
\Pr(Z=0)
=
\Pr(X\leq M)
=
F_X(M).
}
\]

Thus, the reinsurer's payment has a point mass at zero.
\end{frame}

\begin{frame}{Distribution of the Reinsurer's Payment}
\protect\phantomsection\label{distribution-of-the-reinsurers-payment-1}
For \(z>0\),

\[
\begin{aligned}
F_Z(z)
&=\Pr(Z\leq z)\\
&=\Pr((X-M)_+\leq z)\\
&=\Pr(X\leq M+z)\\
&=F_X(M+z).
\end{aligned}
\]

Therefore,

\[
\boxed{
F_Z(z)=
\begin{cases}
0, & z<0,\\
F_X(M), & z=0,\\
F_X(M+z), & z>0.
\end{cases}
}
\]
\end{frame}

\begin{frame}{Continuous Part of the Reinsurer's Payment}
\protect\phantomsection\label{continuous-part-of-the-reinsurers-payment}
For \(z>0\), differentiating gives

\[
\boxed{
f_Z(z)=f_X(M+z).
}
\]

Thus, the distribution of \(Z\) consists of

\[
\boxed{
\text{point mass at }0
+
\text{continuous part for }z>0.
}
\]

The point mass is

\[
\Pr(Z=0)=F_X(M).
\]
\end{frame}

\begin{frame}{Two Different Point Masses}
\protect\phantomsection\label{two-different-point-masses-1}
Under excess-of-loss reinsurance,

\[
Y=\min(X,M)
\]

has a point mass at the retention:

\[
\boxed{
\Pr(Y=M)=1-F_X(M).
}
\]

The reinsurer's payment

\[
Z=(X-M)_+
\]

has a point mass at zero:

\[
\boxed{
\Pr(Z=0)=F_X(M).
}
\]
\end{frame}

\begin{frame}{Interpreting the Point Masses}
\protect\phantomsection\label{interpreting-the-point-masses}
The two point masses have different actuarial interpretations.

\[
Y=M
\]

means that the insurer has reached its retention.

By contrast,

\[
Z=0
\]

means that the claim has not reached the retention.

Both arise from the same excess-of-loss contract.
\end{frame}

\begin{frame}{Example: Distribution of the Reinsurer's Payment}
\protect\phantomsection\label{example-distribution-of-the-reinsurers-payment}
Let

\[
X\sim\operatorname{Exponential}(\lambda)
\]

and

\[
Z=(X-M)_+.
\]

Find:

\begin{enumerate}
\tightlist
\item
  \(F_Z(z)\);
\item
  \(\Pr(Z=0)\);
\item
  the density of the continuous part.
\end{enumerate}
\end{frame}

\begin{frame}{Example: Distribution of the Reinsurer's Payment}
\protect\phantomsection\label{example-distribution-of-the-reinsurers-payment-1}
For \(z>0\),

\[
F_Z(z)
=
F_X(M+z)
=
1-e^{-\lambda(M+z)}.
\]

At zero,

\[
\boxed{
\Pr(Z=0)
=
F_X(M)
=
1-e^{-\lambda M}.
}
\]

For \(z>0\),

\[
\boxed{
f_Z(z)
=
\lambda e^{-\lambda(M+z)}.
}
\]
\end{frame}

\begin{frame}{Mixed Distributions}
\protect\phantomsection\label{mixed-distributions-1}
A random variable can have both:

\begin{itemize}
\tightlist
\item
  a continuous component; and
\item
  one or more point masses.
\end{itemize}

Such a random variable is called a \textbf{mixed distribution}.

For example,

\[
Y=\min(X,M)
\]

has a continuous component below \(M\) and a point mass at \(M\).

Similarly,

\[
Z=(X-M)_+
\]

has a point mass at zero and a continuous component above zero.
\end{frame}

\begin{frame}{Expectations for a Mixed Distribution}
\protect\phantomsection\label{expectations-for-a-mixed-distribution}
Suppose \(U\) has density \(f_U\) for its continuous component and point
masses at \(u_1,u_2,\ldots\).

Then

\[
\boxed{
E[g(U)]
=
\int g(u)f_U(u)\,du
+
\sum_i g(u_i)\Pr(U=u_i).
}
\]

Both the continuous component and the point masses must be included.
\end{frame}

\begin{frame}{Expected Insurer Payment}
\protect\phantomsection\label{expected-insurer-payment}
For

\[
Y=\min(X,M),
\]

the expected payment is

\[
\boxed{
E[Y]
=
\int_0^M x f_X(x)\,dx
+
M\{1-F_X(M)\}.
}
\]

The second term accounts for the point mass at \(M\).
\end{frame}

\begin{frame}
\begin{tcolorbox}[enhanced jigsaw, arc=.35mm, bottomrule=.15mm, bottomtitle=1mm, breakable, colback=white, colbacktitle=quarto-callout-warning-color!10!white, colframe=quarto-callout-warning-color-frame, coltitle=black, left=2mm, leftrule=.75mm, opacityback=0, opacitybacktitle=0.6, rightrule=.15mm, title=\textcolor{quarto-callout-warning-color}{\faExclamationTriangle}\hspace{0.5em}{Do not ignore the point mass}, titlerule=0mm, toprule=.15mm, toptitle=1mm]

It is incorrect to calculate

\[
E[Y]
=
\int_0^M x f_Y(x)\,dx
\]

without accounting for the point mass at \(M\).

The correct expression is

\[
E[Y]
=
\int_0^M x f_X(x)\,dx
+
M\{1-F_X(M)\}.
\]

\end{tcolorbox}
\end{frame}

\begin{frame}{Reinsurance Claims}
\protect\phantomsection\label{reinsurance-claims-1}
The reinsurer's payment

\[
Z=(X-M)_+
\]

is defined for every gross claim.

For claims satisfying \(X\leq M,\) the reinsurer pays zero.

In many actuarial calculations, we are interested specifically in claims
that involve the reinsurer:

\[
\boxed{X>M}.
\]

These are referred to as \textbf{reinsurance claims}.
\end{frame}

\begin{frame}{Reinsurance Payment versus Reinsurance Claim}
\protect\phantomsection\label{reinsurance-payment-versus-reinsurance-claim}
The reinsurer's payment on a randomly selected claim is \(Z=(X-M)_+.\)

This includes the possibility \(Z=0.\)

If we consider only claims for which the reinsurer makes a positive
payment, define

\[
W=Z\mid Z>0.
\]

Since

\[
Z>0
\iff
X>M,
\]

we have

\[
\boxed{
W=X-M\mid X>M.
}
\]
\end{frame}

\begin{frame}{Distribution of a Reinsurance Claim}
\protect\phantomsection\label{distribution-of-a-reinsurance-claim}
For \(w>0\),

\[
\begin{aligned}
F_W(w)
&=\Pr(W\leq w)\\
&=\Pr(X-M\leq w\mid X>M)\\
&=\frac{\Pr(M<X\leq M+w)}
        {\Pr(X>M)}.
\end{aligned}
\]

Therefore,

\[
\boxed{
F_W(w)
=
\frac{F_X(M+w)-F_X(M)}
{1-F_X(M)},
\qquad w>0.
}
\]

This is the distribution of the \textbf{positive excess loss}.
\end{frame}

\begin{frame}{Density of a Reinsurance Claim}
\protect\phantomsection\label{density-of-a-reinsurance-claim}
If \(X\) has density \(f_X\), differentiating gives

\[
\boxed{
f_W(w)
=
\frac{f_X(M+w)}
{1-F_X(M)},
\qquad w>0.
}
\]

Unlike \(Z\), the conditional payment \(W\) has no point mass at zero.

Thus,

\[
\begin{aligned}
Z &: \text{payment on every claim},\\
W &: \text{payment given that the reinsurer is involved}.
\end{aligned}
\]
\end{frame}

\begin{frame}
\begin{tcolorbox}[enhanced jigsaw, arc=.35mm, bottomrule=.15mm, bottomtitle=1mm, breakable, colback=white, colbacktitle=quarto-callout-note-color!10!white, colframe=quarto-callout-note-color-frame, coltitle=black, left=2mm, leftrule=.75mm, opacityback=0, opacitybacktitle=0.6, rightrule=.15mm, title=\textcolor{quarto-callout-note-color}{\faInfo}\hspace{0.5em}{Three quantities to distinguish}, titlerule=0mm, toprule=.15mm, toptitle=1mm]

\[
\begin{aligned}
X &: \text{gross claim amount},\\
Z &: \text{reinsurer's payment on every claim},\\
W &: \text{reinsurer's payment given }X>M.
\end{aligned}
\]

\(Z\) has a point mass at zero, whereas \(W\) describes only positive
reinsurance payments.

\end{tcolorbox}
\end{frame}

\begin{frame}{Example: Distribution of a Reinsurance Claim}
\protect\phantomsection\label{example-distribution-of-a-reinsurance-claim}
Suppose \(X\) has distribution function \(F_X\) and density \(f_X\). The
insurer purchases excess-of-loss reinsurance with retention \(M\).

Let \(W=Z\mid Z>0.\)

Show that

\[
F_W(w)
=
\frac{F_X(M+w)-F_X(M)}
{1-F_X(M)},
\qquad w>0,
\]

and

\[
f_W(w)
=
\frac{f_X(M+w)}
{1-F_X(M)},
\qquad w>0.
\]

Interpret \(W\) as the amount by which a claim exceeds the insurer's
retention, given that the retention has been exceeded.
\end{frame}

\begin{frame}{Expected Payments}
\protect\phantomsection\label{expected-payments-1}
Under excess-of-loss reinsurance,

\[
Y=\min(X,M)  \quad \text{ and } \quad   Z=(X-M)_+.
\]

Since \(X=Y+Z,\) taking expectations gives

\[
\boxed{
E[X]=E[Y]+E[Z].
}
\]

Hence, \(\quad \quad \boxed{E[Y]=E[X]-E[Z].}\)
\end{frame}

\begin{frame}{Expected Reinsurer Payment}
\protect\phantomsection\label{expected-reinsurer-payment}
The reinsurer pays

\[
Z=(X-M)_+.
\]

Therefore,

\[
\boxed{
E[Z]
=
\int_M^\infty
(x-M)f_X(x)\,dx.
}
\]

Using the substitution \(w=x-M\),

\[
\boxed{
E[Z]
=
\int_0^\infty
w f_X(M+w)\,dw.
}
\]

This is the expected excess loss above the retention \(M\).
\end{frame}

\begin{frame}{Expected Insurer Payment}
\protect\phantomsection\label{expected-insurer-payment-1}
The insurer pays \(Y=\min(X,M).\)

Hence,

\[
\boxed{
E[Y]
=
\int_0^M x f_X(x)\,dx
+
M\{1-F_X(M)\}.
}
\]

Alternatively, since \(Y=X-Z,\) we have

\[
\boxed{
E[Y]=E[X]-E[Z].
}
\]
\end{frame}

\begin{frame}
\begin{tcolorbox}[enhanced jigsaw, arc=.35mm, bottomrule=.15mm, bottomtitle=1mm, breakable, colback=white, colbacktitle=quarto-callout-note-color!10!white, colframe=quarto-callout-note-color-frame, coltitle=black, left=2mm, leftrule=.75mm, opacityback=0, opacitybacktitle=0.6, rightrule=.15mm, title=\textcolor{quarto-callout-note-color}{\faInfo}\hspace{0.5em}{Interpretation of the expected payments}, titlerule=0mm, toprule=.15mm, toptitle=1mm]

Under excess-of-loss reinsurance,

\[
E[Z]
\]

is the expected amount transferred to the reinsurer per gross claim.

Similarly,

\[
E[Y]
\]

is the expected amount retained by the direct insurer.

Therefore,

\[
\boxed{
E[X]=E[Y]+E[Z].
}
\]

Reinsurance changes how the loss is shared between the insurer and
reinsurer.

\end{tcolorbox}
\end{frame}

\begin{frame}{Example: Expected Insurer and Reinsurer Payments}
\protect\phantomsection\label{example-expected-insurer-and-reinsurer-payments}
Let \(X\) have density \(f_X\), and suppose the insurer purchases
unlimited excess-of-loss reinsurance with retention \(M\).

Show that

\[
E[Z]
=
\int_M^\infty
(x-M)f_X(x)\,dx
\]

and

\[
E[Y]
=
E[X]-E[Z].
\]

Then verify that \(E[X]=E[Y]+E[Z].\)
\end{frame}

\begin{frame}{Exponential Claims}
\protect\phantomsection\label{exponential-claims-1}
Suppose

\[
X\sim\operatorname{Exponential}(\lambda).
\]

Then

\[
\boxed{
E[X]=\frac{1}{\lambda}
}
 \quad \text{ and } \quad 
\boxed{
\Pr(X>M)=e^{-\lambda M}.
}
\]

Thus, the probability that a claim involves the reinsurer is

\[
\boxed{
\Pr(X>M)=e^{-\lambda M}.
}
\]
\end{frame}

\begin{frame}{Expected Reinsurer Payment}
\protect\phantomsection\label{expected-reinsurer-payment-1}
For exponential claims,

\[
E[Z]
=
\int_M^\infty
(x-M)\lambda e^{-\lambda x}\,dx.
\]

This gives

\[
\boxed{
E[Z]
=
\frac{e^{-\lambda M}}{\lambda}.
}
\]

Since \(E[X]=\frac{1}{\lambda},\)

we can write

\[
\boxed{
E[Z]=e^{-\lambda M}E[X].
}
\]
\end{frame}

\begin{frame}{Expected Insurer Payment}
\protect\phantomsection\label{expected-insurer-payment-2}
Using

\[
E[Y]=E[X]-E[Z],
\]

we obtain

\[
\boxed{
E[Y]
=
\frac{1-e^{-\lambda M}}{\lambda}.
}
\]

Thus,

\[
\boxed{
E[X]
=
\frac{1-e^{-\lambda M}}{\lambda}
+
\frac{e^{-\lambda M}}{\lambda}.
}
\]

The expected gross claim is divided between the insurer and reinsurer.
\end{frame}

\begin{frame}{Interpreting the Retention Level}
\protect\phantomsection\label{interpreting-the-retention-level}
As \(M\) increases, \(\Pr(X>M)=e^{-\lambda M}\) decreases.

Therefore:

\begin{itemize}
\tightlist
\item
  fewer claims involve the reinsurer; the expected reinsurer payment
  decreases; and
\item
  the expected insurer payment increases.
\end{itemize}

The retention \(M\) controls the amount of risk transferred to the
reinsurer.

For exponential claims,

\[
E[Z]
=
\frac{e^{-\lambda M}}{\lambda}.
\]

The simple form follows from the memoryless property of the exponential
distribution.
\end{frame}




\end{document}
